% Generated draft: requires mathematical and source review.
\documentclass[11pt]{article}
\usepackage[margin=25mm]{geometry}
\usepackage{fontspec}
\setmainfont{Latin Modern Roman}
\usepackage{amsmath,amssymb,mathtools,unicode-math}
\setmathfont{Latin Modern Math}
\usepackage{enumitem,needspace,xurl,hyperref}
\setlength{\parindent}{0pt}
\setlength{\parskip}{5pt}
\setlength{\emergencystretch}{3em}
\newcommand{\N}{\mathbb{N}}
\newcommand{\Z}{\mathbb{Z}}
\newcommand{\Q}{\mathbb{Q}}
\newcommand{\R}{\mathbb{R}}
\newcommand{\C}{\mathbb{C}}
\newcommand{\F}{\mathbb{F}}
\newcommand{\A}{\mathbb{A}}
\newcommand{\PP}{\mathbb{P}}
\newcommand{\E}{\mathbb{E}}
\newcommand{\eps}{\varepsilon}
\newcommand{\dd}{\,\mathrm{d}}
\begin{document}
\section*{UnsolvedMath ID: 30001606}
\textbf{Source problem code:} catalogue entry
\subsection{Definitions and mathematical statement}
Let \(n\geq 1\) be an integer and fix a vector \(\ell=(\ell_1,\ldots,\ell_n)\in(0,\infty)^n\). Write \(S^2=\{v\in\R^3:\|v\|=1\}\), and define
\[
 P_\ell=\left\{(v_1,\ldots,v_n)\in(S^2)^n:\sum_{i=1}^n\ell_i v_i=0\right\},\qquad N_\ell=P_\ell/SO(3).
\]
Here \(SO(3)\) acts diagonally by \(g\cdot(v_1,\ldots,v_n)=(gv_1,\ldots,gv_n)\), and \(N_\ell\) carries the quotient topology. Thus configurations differing by a common rotation represent the same point.

To specify the numerical convention for higher topological complexity, let \(X\) be a nonempty path-connected space and let \(r\geq2\) be an integer. Give \(PX=C([0,1],X)\) the compact-open topology and define
\[
 e_r:PX\longrightarrow X^r,\qquad e_r(\gamma)=\left(\gamma(0),\gamma\!\left(\frac1{r-1}\right),\ldots,\gamma(1)\right).
\]
The unreduced sequential topological complexity \(\operatorname{TC}^{\mathrm{unred}}_r(X)\) is the least positive integer \(k\) for which \(X^r\) has an open cover \(U_1,\ldots,U_k\), each admitting a continuous map \(s_j:U_j\to PX\) satisfying \(e_r\circ s_j=\operatorname{id}_{U_j}\). Set this value to \(\infty\) if no finite such cover exists. The reduced convention counts one fewer open set: for finite values,
\[
 \operatorname{TC}^{\mathrm{red}}_r(X)=\operatorname{TC}^{\mathrm{unred}}_r(X)-1.
\]

\textbf{Problem as supplied.} For the fixed positive length vector \(\ell\), compute \(\operatorname{TC}_n(N_\ell)\).

\textbf{Unresolved qualifications.} The source record does not specify reduced versus unreduced normalization. It also uses \(n\) both for the number of lengths and for the sequential-complexity index. Read literally, the requested index equals the number of lengths; whether the intended question instead asks for \(\operatorname{TC}_r(N_\ell)\) for every \(r\geq2\), independently of \(n\), remains unresolved. No lower bound \(n\geq2\), genericity hypothesis on \(\ell\), nonemptiness hypothesis, or path-connectedness hypothesis is supplied. Consequently, the definition above specifies the invariant when its stated hypotheses hold, but does not silently impose those hypotheses on the source problem or settle conventions for empty spaces or index \(1\). Degenerate length vectors are not excluded by the supplied formulation.
\subsection{Short English statement}
For a fixed vector of positive side lengths, form the space of closed spatial polygon configurations modulo simultaneous rotations and compute its higher topological complexity. The draft preserves the source's request \(\operatorname{TC}_n(N_\ell)\), while flagging the unspecified normalization, possible index ambiguity, and missing hypotheses on the polygon space. The supplied record labels the problem partially solved; this is a dated curation assessment requiring independent review.
\subsection{Sources}
\begin{itemize}
\item Local UnsolvedMath registry, supplied record with permanent ID \(30001606\), source problem code \texttt{catalogue entry}, in the collection \emph{Oberwolfach Reports Open Problems}; title: \emph{Higher Topological Complexity of Spatial Polygon Spaces}.
\item Supplied primary citation: \emph{Arbeitsgemeinschaft: Topological Robotics} (2011), \emph{Oberwolfach Reports}, pp.\ 2681--2744, DOI \texttt{10.4171/owr/2010/47}: \url{https://doi.org/10.4171/owr/2010/47}. The background describes this collection as workshop reports documenting open-problem discussions. The record gives proposed year 2010 and names no individual proposer.
\item Additional supplied literature link: \url{https://doi.org/10.1016/j.topol.2017.09.022}. No bibliographic title or authors are supplied for this link.
\item Additional supplied literature link: \url{https://doi.org/10.1007/s41468-026-00231-5}. No bibliographic title or authors are supplied for this link.
\item Literature assessment dated exactly \texttt{2026-08-21}: the supplied curation labels the problem partially solved and states that higher sequential topological complexity has been computed for many spaces, that ordinary topological complexity of nonorientable surfaces is known in most genera, and that the full family denoted \(\operatorname{TC}_n(N_l)\) is not determined. This is supplied dated curation, not an independently verified literature assessment. The supplied background itself requests independent verification of time-sensitive status and cited claims.
\item The supplied statement status is \texttt{corrected\_verified}; its verification note says the mathematical content was checked against the official Oberwolfach report and the notation, context, and question form were normalized without changing the source-level claim. This is a report of supplied verification metadata, not a claim that the present draft independently consulted the report.
\end{itemize}
\end{document}
